%%
% BIThesis Undergraduate Thesis Template (English) — Compile with XeLaTeX
% Chapter 1: Introduction
%%

\chapter{Introduction}

\section{The IXP Bandwidth Resale Problem}

Internet Exchange Points (IXPs) are critical infrastructure where networks peer to exchange traffic directly \cite{lantzNetworkSoftwareMust2010, martinIXPPeeringAgreements2014}. Large enterprises and content delivery networks operate massive networks with global presence, while smaller networks and emerging providers lack the scale to negotiate peering directly with these giants. The traditional solution is to purchase ``transit'' bandwidth---a small network contracts with a transit provider for a fixed committed amount of bandwidth, such as 1 Gbps, under a long-term agreement.

However, this model introduces significant inefficiency. Bandwidth demand varies throughout the day: a small network might need 1 Gbps during peak hours but only 200 Mbps during off-peak periods, resulting in 80\% capacity waste during low-demand windows. Conversely, unexpected traffic spikes can exceed contracted capacity, forcing expensive emergency purchases or resulting in severe packet loss. This demand volatility, driven by diurnal traffic patterns and unpredictable network events, creates a highly lucrative market opportunity for dynamic bandwidth resale.

A transit reseller purchases bulk bandwidth from major providers (e.g., 100 Gbps from major carriers) and subdivides it among many smaller customers. If this subdivision is static, the reseller faces the same inefficiencies as the customers. To maximize link utilization and profitability, the reseller requires a dynamic, competitive allocation mechanism. This thesis addresses this need by designing a Software-Defined Networking (SDN) global control plane that treats network bandwidth as a dynamic commodity, allocated via periodic economic auctions.

\section{SDN and Microservice Architectures}

To facilitate a global, high-throughput bandwidth marketplace, the control plane must be highly scalable and resilient. Traditional monolithic SDN controllers \cite{mckeownOpenFlowEnablingInnovation2008} create single points of failure and struggle to scale individual logical components (such as bid ingestion versus flow metric processing). 

Modern cloud-native design dictates that the control plane be decomposed into isolated microservices \cite{newmanBuildingMicroservices2015}. By utilizing a state-separated architecture---where stateless microservices communicate via an asynchronous message broker (Apache Kafka) \cite{newtmannKafkaDistributedStreamingPlatform2016} and rely on a distributed consensus store (Atomix) \cite{ongaroRaftConsensusAlgorithm2014} for a single source of truth---the system achieves both horizontal scalability and strict fault isolation. 

\section{The Observability Gap in Asynchronous Systems}

While asynchronous microservices provide superior scalability, they introduce severe operational blind spots. In a traditional synchronous system, a failing component crashes, and standard infrastructure orchestrators (like Kubernetes) simply restart it. 

However, state-separated SDN controllers frequently suffer from ``soft failures.'' For example, if the Atomix consensus database experiences severe network latency, the surrounding microservices do not crash; they simply hang, waiting for database locks. Standard infrastructure monitoring tools \cite{fernandoContainerOrchestrationKubernetes2017} (e.g., CPU/Memory metrics and \texttt{LivenessProbes}) report the system as perfectly healthy, while the actual business logic (the auction clearing) is silently paralyzed.

Furthermore, because communication is asynchronous, traditional request tracing fails \cite{sigelman2010dapper}. When an API Gateway writes a bid to a database, and an Auction Runner reads it 60 seconds later, the causal link between the customer's HTTP request and the auction algorithm's execution is lost. Bridging this observability gap in asynchronous environments is the primary technical challenge addressed by this thesis \cite{majorsObservabilityEngineering2022}.

\section{Terminology Definitions}

Because this thesis spans both network engineering and distributed software engineering, it is critical to explicitly differentiate between two distinct types of telemetry utilized within the system architecture:

\begin{itemize}
    \item \textbf{Network Telemetry (Business Logic):} Refers to data-plane traffic statistics (e.g., bytes transferred, flow rates in Kbps) periodically emitted by the physical or simulated switches. This data is processed by the control plane and served to customer agents to inform their economic bidding strategies.
    \item \textbf{Observability Telemetry (System Health):} Refers to the operational signals (traces, logs, and infrastructure metrics) generated by the OpenTelemetry SDKs embedded within the Go microservices. This data is used exclusively by system operators to monitor the software's health, debug latency, and detect application-layer failures.
\end{itemize}

\section{Objectives and Contributions}

The primary objective of this thesis is to design, implement, and evaluate a microservice-based SDN global control plane that is both economically efficient and fully observable. The core contributions are as follows:

\begin{enumerate}
    \item \textbf{State-Separated Auction Control Plane:} The design and implementation of a microservice architecture that dynamically shapes network traffic based on a periodic, sealed-bid uniform-price auction algorithm.
    \item \textbf{Asynchronous Trace Context Propagation:} The development of a custom W3C TraceContext propagation strategy that utilizes the Atomix consensus database as a context carrier, enabling unbroken distributed tracing across temporally decoupled services.
    \item \textbf{Observability-Driven Fault Localization:} The integration of a comprehensive OpenTelemetry framework (Metrics, Traces, and Logs) to explicitly detect and localize soft failures (consensus degradation) that evade traditional container orchestration monitoring.
\end{enumerate}

\section{Thesis Organization}

The remainder of this thesis is structured as follows:
\begin{itemize}
    \item \textbf{Chapter 2 (System Architecture):} Details the high-level microservice topology, the asynchronous decoupling strategy using Apache Kafka, and the primary data workflows.
    \item \textbf{Chapter 3 (System Implementation):} Provides an in-depth analysis of the control plane's core logic, including the mathematical formulation of the uniform-price auction and the custom trace-linking code.
    \item \textbf{Chapter 4 (Observability Framework Design):} Presents the OpenTelemetry architecture, detailing the data reduction strategies (tail-sampling and severity-filtering) used to minimize performance overhead.
    \item \textbf{Chapter 5 (System Evaluation):} Evaluates the system's functional correctness during a massive traffic spike, and validates the observability framework's ability to detect an injected database latency fault.
    \item \textbf{Chapter 6 (Conclusions):} Summarizes the contributions and proposes directions for future research.
\end{itemize}